\chapter*{Введение}
\addcontentsline{toc}{chapter}{Введение}

Современная биоинформатика работает с данными огромного масштаба. Геномы, транскриптомы и белковые последовательности становятся все доступнее, однако способы их восприятия остаются преимущественно визуальными и числовыми. Для специалиста последовательность из миллионов символов является материалом для вычислительного анализа, требующим сжатия, сравнения и статистической обработки. В связи с этим возникает задача поиска дополнительных способов представления биологических данных, которые позволяли бы воспринимать их структуру не только через таблицы, графики и численные признаки.

Музыкальная сонификация предлагает один из таких способов. Человек способен воспринимать сложные временные структуры слухом: высоту, длительность, повтор, плотность, динамику и изменение регистра. Поэтому биологическую последовательность можно рассматривать не только как строку символов, но и как источник управляющих признаков для построения музыкальной формы. Актуальность данной работы определяется ростом объема биологических данных, интересом к сонификации как средству научной коммуникации и развитием генеративных моделей, которые позволяют уходить от прямого отображения символов в ноты к более осмысленной музыкальной генерации.

Задача сонификации биологических данных исследуется с 1980-х годов. Ранние работы использовали прямое отображение нуклеотидов в ноты~\cite{ohno1986}. Современные подходы включают алгоритмическую композицию~\cite{temple2017}, параметрическую сонификацию~\cite{grond2011} и генеративные модели. В области генерации символической музыки активно развиваются методы на основе рекуррентных сетей~\cite{eck2002}, вариационных автоэнкодеров~\cite{roberts2018}, генеративно-состязательных сетей~\cite{dong2018} и трансформеров~\cite{huang2018}. Обзорные работы по deep music generation показывают, что качество результата зависит не только от мощности модели, но и от музыкального представления, метрик качества и устойчивости обучения~\cite{briot2017,ji2020}. Датасеты POP909~\cite{wang2020} и MAESTRO~\cite{hawthorne2019} стали распространенной основой для обучения моделей символической музыки.

На предварительном этапе работы рассматривались вопросы выбора музыкальной нотации и оценки результата: MusicLang как высокоуровневый язык описания музыки, MIDI как универсальный событийный формат, а также идея модели-оценщика для музыкальных фрагментов. Проведенный анализ показал, что для финальной системы важнее не прямое кодирование символов биологической последовательности, а промежуточное структурированное представление, пригодное для обучения и последующей количественной оценки. Существующие решения часто используют слишком простое отображение символов или не учитывают музыкальную структуру, поэтому задача создания системы, которая генерирует музыкально связные композиции под управлением биологических признаков, остается актуальной.

Объектом исследования являются биологические последовательности и способы их преобразования в воспринимаемую человеком форму. Предметом исследования является метод нейросетевой генерации символической музыки, условно управляемой признаками биологических последовательностей.

Цель работы состоит в разработке и экспериментальной проверке обучаемой нейросетевой модели для синтеза и оценки мелодий по биологическим данным, в которой биологические признаки используются как управляющий сигнал для генерации музыки. Для достижения этой цели необходимо изучить предметную область биосонификации и генерации символической музыки, проанализировать существующие решения и выявить их ограничения, подготовить биологические данные в формате FASTA и музыкальный корпус MIDI, реализовать извлечение признаков из DNA, RNA и protein последовательностей, разработать иерархическую архитектуру генерации Bio->Harmony->Melody, реализовать сопоставление биологических и музыкальных сегментов, обучить модели на доступном оборудовании, провести экспериментальную оценку и сравнение с baseline, а также подготовить воспроизводимые результаты и документацию.

В утвержденной теме работы используется термин ``самообучающаяся модель''. На начальном этапе проектирования рассматривалась генеративно-состязательная схема, в которой одна модель генерирует музыкальный материал, а другая модель-оценщик формирует обучающий сигнал для улучшения генератора. Такая постановка соответствует идее итеративного улучшения за счет внутренней оценки результата. В ходе анализа предметной области и вычислительных ограничений итоговая архитектура была изменена: для задачи генерации символической музыки по биологическим данным была выбрана более устойчивая и воспроизводимая схема на основе условного autoregressive Transformer, обучаемого на парах ``биологический фрагмент - музыкальный сегмент''. Поэтому далее термин ``самообучающаяся'' трактуется не как автономное онлайн-обучение в процессе работы веб-интерфейса и не как финальная GAN-архитектура, а как обучаемая на данных нейросетевая модель, параметры которой автоматически оптимизируются по функции потерь на обучающей выборке. Оценка качества вынесена в экспериментальный контур и проводится по заранее заданным метрикам.

В работе используются анализ научной литературы и существующих решений, проектирование программной архитектуры, реализация программного комплекса на Python, методы машинного обучения с использованием PyTorch, экспериментальная оценка на тестовых данных и сравнительный анализ с baseline.

Научная новизна работы состоит в предложении иерархической архитектуры, где биологические признаки используются как управляющий сигнал для структурированной музыкальной генерации. В отличие от прямого отображения символов, система разделяет задачу на два этапа: генерацию гармонии и генерацию мелодии поверх гармонии. Это обеспечивает музыкальную связность при сохранении управления от биологических данных.

Практическая значимость работы состоит в создании рабочего программного прототипа BioSonification. Система включает обученные модели с воспроизводимыми checkpoint, веб-интерфейс для генерации музыки из FASTA, документацию и примеры использования, а также отчеты по данным и метрикам качества. Система может использоваться в образовательных целях, для научной коммуникации и демонстрации биоинформатических данных. Описание программы, руководство пользователя и руководство разработчика приведены в приложении А.

Работа состоит из введения, трех глав, заключения, списка использованных источников и литературы и приложения. В первой главе анализируются предметная область, существующие подходы к биосонификации и генерации музыки, выявляются недостатки аналогов и формулируется постановка задачи. Во второй главе описываются проектирование и реализация системы: техническое задание, архитектура, биологический энкодер, музыкальное представление, pairing, нейросетевая модель и процесс обучения. В третьей главе представлена экспериментальная оценка: методика проверки, метрики, результаты обучения, сравнение с baseline и анализ ограничений. В заключении подводятся итоги работы и намечаются направления дальнейшего развития.
